% !TEX program = pdflatex
% ════════════════════════════════════════════════════════════════════════════
%  UrbanFlow CDMX / VialAI — Interfaz de uso
%  Diplomado en Ciencia de Datos · FES Acatlán, UNAM · Gen 30 · 2026
%  Autor: Carlos Armando López Encino
%  Profesor: Mtro. Fernando Barranco Rodríguez
% ════════════════════════════════════════════════════════════════════════════

\documentclass[11pt,a4paper]{article}

% ───── Codificación y lenguaje ─────────────────────────────────
\usepackage[utf8]{inputenc}
\usepackage[T1]{fontenc}
\usepackage[spanish,es-tabla,es-noindentfirst]{babel}
\usepackage{lmodern}
\usepackage{microtype}

% ───── Matemáticas ─────────────────────────────────────────────
\usepackage{amsmath}
\usepackage{siunitx}
\sisetup{output-decimal-marker={.},group-separator={\,}}

% ───── Geometría y espaciado ───────────────────────────────────
\usepackage[a4paper,margin=2.5cm,headheight=15pt]{geometry}
\usepackage{setspace}
\setstretch{1.15}
\usepackage{parskip}

% ───── Color e identidad VialAI ────────────────────────────────
\usepackage{xcolor}
\definecolor{vialNavyDeep}{HTML}{0B1F3A}
\definecolor{vialNavy}{HTML}{0F2E4D}
\definecolor{vialTeal}{HTML}{2EBFA6}
\definecolor{vialAqua}{HTML}{3BA3BD}
\definecolor{vialGreen}{HTML}{A3D88F}
\definecolor{vialAmber}{HTML}{F5B836}
\definecolor{vialText}{HTML}{1F2937}
\definecolor{vialMuted}{HTML}{4A5B6D}
\definecolor{vialBgSoft}{HTML}{F3F6FA}

% ───── Gráficos ────────────────────────────────────────────────
\usepackage{graphicx}
\graphicspath{{./}{imagenes_12/}{imagenes_12/imagenes_png/}}
\usepackage{float}

% ───── Tablas ──────────────────────────────────────────────────
\usepackage{booktabs}
\usepackage{array}
\usepackage{tabularx}
\usepackage{multirow}
\usepackage{colortbl}

% ───── Tipografía de secciones ─────────────────────────────────
\usepackage{titlesec}
\titleformat{\section}
  {\color{vialNavy}\normalfont\Large\bfseries}{\thesection.}{0.5em}{}
\titleformat{\subsection}
  {\color{vialNavy}\normalfont\large\bfseries}{\thesubsection}{0.5em}{}
\titleformat{\subsubsection}
  {\color{vialTeal}\normalfont\normalsize\bfseries}{\thesubsubsection}{0.5em}{}

% ───── Encabezados y pies ──────────────────────────────────────
\usepackage{fancyhdr}
\pagestyle{fancy}
\fancyhf{}
\fancyhead[L]{\color{vialMuted}\small VialAI --- Interfaz de uso}
\fancyhead[R]{\color{vialMuted}\small\thepage}
\renewcommand{\headrulewidth}{0.4pt}
\renewcommand{\headrule}{\color{vialTeal}\hrule width\headwidth}

% ───── Código ──────────────────────────────────────────────────
\usepackage{listings}
\lstdefinestyle{vialcode}{
  backgroundcolor=\color{vialBgSoft},
  basicstyle=\ttfamily\footnotesize\color{vialText},
  keywordstyle=\color{vialNavy}\bfseries,
  stringstyle=\color{vialTeal},
  commentstyle=\color{vialMuted}\itshape,
  frame=single, framesep=5pt, rulecolor=\color{vialTeal},
  numbers=left, numberstyle=\tiny\color{vialMuted},
  breaklines=true, showstringspaces=false, captionpos=b,
}
\lstset{style=vialcode,
  literate={á}{{\'a}}1 {é}{{\'e}}1 {í}{{\'i}}1 {ó}{{\'o}}1 {ú}{{\'u}}1
           {Á}{{\'A}}1 {É}{{\'E}}1 {Í}{{\'I}}1 {Ó}{{\'O}}1 {Ú}{{\'U}}1
           {ñ}{{\~n}}1 {Ñ}{{\~N}}1 {¿}{?`}1 {¡}{!`}1
}

% ───── Cajas de callout ────────────────────────────────────────
\usepackage[most]{tcolorbox}

\newtcolorbox{tipbox}[1][]{%
  enhanced, breakable,
  colback=vialBgSoft, colframe=vialTeal, coltitle=vialNavy,
  fonttitle=\bfseries, title=#1,
  boxrule=0.8pt, arc=3mm, left=4mm, right=4mm, top=2mm, bottom=2mm,
}

\newtcolorbox{keybox}[1][]{%
  enhanced, breakable,
  colback=vialNavy!5, colframe=vialNavy, coltitle=white,
  fonttitle=\bfseries, title=#1,
  boxrule=1pt, arc=3mm, left=4mm, right=4mm, top=2mm, bottom=2mm,
}

\newtcolorbox{casebox}[1][]{%
  enhanced, breakable,
  colback=vialAmber!10, colframe=vialAmber!80!black, coltitle=vialNavy,
  fonttitle=\bfseries, title=#1,
  boxrule=0.8pt, arc=3mm, left=4mm, right=4mm, top=2mm, bottom=2mm,
}

% ───── Hipervínculos ───────────────────────────────────────────
\usepackage[colorlinks=true,
            linkcolor=vialNavy,
            citecolor=vialTeal,
            urlcolor=vialAqua]{hyperref}

% ───── Comandos propios ────────────────────────────────────────
\newcommand{\VialAI}{\textbf{\textcolor{vialTeal}{Vial}\textcolor{vialAmber}{AI}}}
\newcommand{\ZMVM}{\textsc{zmvm}}
\newcommand{\CDMX}{\textsc{cdmx}}

% ──────────────────────────────────────────────────────────────
\begin{document}

% ═══════════════════════ PORTADA ═══════════════════════════════
\begin{titlepage}
\centering
\vspace*{1.5cm}

{\color{vialMuted}\large Diplomado en Ciencia de Datos \textbullet\ Generación 30}\\[0.3em]
{\color{vialMuted}\small FES Acatlán \textbullet\ UNAM \textbullet\ 2026}\\[2cm]

{\color{vialTeal}\Huge\bfseries UrbanFlow \CDMX\ / \VialAI}\\[0.8em]
{\color{vialNavy}\LARGE\bfseries Interfaz de uso}\\[0.5em]
{\color{vialMuted}\large\itshape Cómo se construyó el sistema y\\
cómo el usuario lo opera}\\[2.5cm]

{\large\textbf{Carlos Armando López Encino}}\\
{\small\color{vialMuted}cale.delta@gmail.com}\\[0.6em]
{\small Profesor: Mtro. Fernando Barranco Rodríguez}\\[2cm]

{\color{vialMuted}\small Proyecto integrador \textbullet\ 18 de abril de 2026}\\[0.4em]
{\color{vialAqua}\small\url{https://github.com/caledelta/UrbanFlow_CDMX}}

\vfill
{\color{vialMuted}\footnotesize
Documento complementario al artículo técnico \texttt{proyecto\_ZMCDMX\_v3.pdf}.
Cubre el criterio ``Interfaz de uso'' de la rúbrica de evaluación.
Para ejecutar el sistema: \texttt{streamlit run app/streamlit\_app.py}}
\end{titlepage}

\newpage
\tableofcontents
\newpage

% ═══════════════════════════════════════════════════════════════
%  PARTE I: CÓMO SE CONSTRUYÓ
% ═══════════════════════════════════════════════════════════════

\part*{Parte I — Cómo se construyó el sistema}
\addcontentsline{toc}{part}{Parte I: Cómo se construyó el sistema}

% ═══════════════════════ §1 VISIÓN GENERAL ═════════════════════
\section{Visión general de la arquitectura}

\VialAI\ es un sistema de cinco capas que transforma datos de tráfico
en tiempo real en una recomendación conversacional:

\begin{keybox}[Las cinco capas del sistema]
\begin{enumerate}
\item \textbf{Ingesta de datos} (\texttt{src/ingestion/}) ---
cinco fuentes externas: TomTom Traffic Stats, TomTom Routing,
OpenWeatherMap, C5 \CDMX, Google Maps Distance Matrix.
\item \textbf{Motor estocástico} (\texttt{src/simulation/}) ---
cadena de Markov de 3 estados (Fluido / Lento / Congestionado)
calibrada con MLE sobre \num{1}M+ registros del C5; simulación
Monte Carlo vectorizada de \num{10000} trayectorias.
\item \textbf{Perturbaciones contextuales} (\texttt{src/agent/perturbaciones.py}) ---
catálogo de 11 eventos típicos de la \ZMVM\ (marchas, cierres de
Metro, festivos patrios) que modifican multiplicativamente la
matriz de transición.
\item \textbf{Agente conversacional} (\texttt{src/agent/agent.py}) ---
bucle de \emph{tool use} sobre Claude claude-sonnet-4-6 con 6 herramientas
especializadas, inyección de contexto de lugares personalizados y
razonamiento sobre los resultados del motor.
\item \textbf{Interfaz multimodal} (\texttt{app/streamlit\_app.py}) ---
dashboard Streamlit con mapa Folium interactivo, tres modos de
entrada (manual, chat, voz) y telemetría local consentida.
\end{enumerate}
\end{keybox}

\begin{table}[H]
\centering
\small
\renewcommand{\arraystretch}{1.4}
\begin{tabularx}{\linewidth}{@{}l l X@{}}
\toprule
\rowcolor{vialNavy!10}
\textbf{Capa} & \textbf{Tecnología principal} & \textbf{Módulo} \\
\midrule
Backend numérico & Python 3.14, NumPy, SciPy & \texttt{src/simulation/} \\
Datos tabulares & Pandas & \texttt{src/ingestion/} \\
Validación de esquemas & Pydantic 2.x & \texttt{src/models/schemas.py} \\
Agente LLM & Anthropic SDK (Claude claude-sonnet-4-6) & \texttt{src/agent/agent.py} \\
Voz STT & faster-whisper (local) & \texttt{src/agent/voice\_io.py} \\
Voz TTS & pyttsx3 (local) / OpenAI nova (premium) & \texttt{src/agent/voice\_io.py} \\
Frontend & Streamlit 1.56 & \texttt{app/streamlit\_app.py} \\
Mapas interactivos & Folium & \texttt{app/streamlit\_app.py} \\
Visualizaciones & Plotly & \texttt{app/streamlit\_app.py} \\
Testing & pytest (\num{720} tests) & \texttt{tests/} \\
\bottomrule
\end{tabularx}
\caption{Stack tecnológico de \VialAI\ por capa funcional.}
\label{tab:stack}
\end{table}

% ═══════════════════════ §2 PROCESO DE CONSTRUCCIÓN ════════════
\section{Proceso de construcción: 16 bloques incrementales}

El sistema se construyó en \textbf{16 bloques de desarrollo
incrementales}, cada uno con su propio commit en GitHub y con
verificación manual de que los tests seguían en verde antes de
avanzar. Este proceso garantiza que el historial del repositorio
es un registro auditable de cada decisión de diseño.

\begin{table}[H]
\centering
\small
\renewcommand{\arraystretch}{1.4}
\begin{tabularx}{\linewidth}{@{}r X l@{}}
\toprule
\rowcolor{vialNavy!10}
\textbf{Bloque} & \textbf{Qué se construyó} & \textbf{Capa} \\
\midrule
1--3  & Motor Monte Carlo + cadena de Markov + calibración MLE & Simulación \\
4--5  & Conectores TomTom (Traffic + Routing) + pipeline ETL & Ingesta \\
6     & Perturbaciones contextuales (catálogo de 11 eventos) & Agente \\
7--8  & Agente VialAI con tool-use loop (6 herramientas) & Agente \\
9     & Dashboard Streamlit inicial + mapa Folium & Interfaz \\
10--12 & Interfaz de voz de tres capas (STT + TTS con fallback) & Voz \\
13    & Rutas personalizadas y lugares guardados & Core \\
14--15 & Sistema de recompensa silenciosa + telemetría + iconos & Core \\
16    & Fixes de UX (flujo Iniciar viaje / Llegué) + marcadores & Interfaz \\
\bottomrule
\end{tabularx}
\caption{Los 16 bloques de desarrollo de \VialAI. Cada bloque corresponde a uno
         o más commits en \url{https://github.com/caledelta/UrbanFlow_CDMX}.}
\label{tab:bloques}
\end{table}

\begin{tipbox}[Principio de desarrollo aplicado]
Commits pequeños, pruebas automáticas entre cada paso, rollback
inmediato ante cualquier regresión. Los \num{720} tests en
\texttt{tests/} cubren el motor estocástico, el pipeline de datos
y el agente conversacional. La velocidad de desarrollo sostenida
se logra \emph{no} con grandes saltos, sino con iteraciones pequeñas
y auditables.
\end{tipbox}

% ═══════════════════════ §3 MOTOR ESTOCÁSTICO ══════════════════
\section{El motor estocástico: Markov + Monte Carlo}

\subsection{Cadena de Markov de tráfico}

El tráfico se modela como una cadena de Markov discreta de tres
estados $S \in \{0 = \text{Fluido}, 1 = \text{Lento}, 2 = \text{Congestionado}\}$.
La matriz de transición $P$ se calibra por máxima verosimilitud
sobre la base histórica del C5 \CDMX\ (\num{1}M+ registros,
2018--2024), segmentada por corredor, día de semana y franja horaria.

\begin{table}[H]
\centering
\small
\renewcommand{\arraystretch}{1.4}
\begin{tabular}{@{}l c c c@{}}
\toprule
\rowcolor{vialNavy!10}
\textbf{Estado actual} $\backslash$ \textbf{Estado siguiente}
  & \textbf{Fluido} & \textbf{Lento} & \textbf{Congestionado} \\
\midrule
Fluido       & 0.72 & 0.21 & 0.07 \\
Lento        & 0.18 & 0.52 & 0.30 \\
Congestionado & 0.05 & 0.35 & 0.60 \\
\bottomrule
\end{tabular}
\caption{Matriz de transición base $P$ para día hábil, pico AM (valores ilustrativos
         a partir de calibración MLE sobre datos C5 \CDMX).}
\label{tab:markov}
\end{table}

\subsection{Perturbaciones contextuales}

Cuando se activa un evento (por ejemplo, ``Marcha 9 de marzo''),
se aplica un factor multiplicador $f$ sobre $P[i, 2]$ (probabilidad
de transitar a Congestionado) y se renormalizan las filas:

\[
P'[i, 2] = \min(P[i, 2] \times f,\; 1), \qquad
P'[i, 0\ldots1] \text{ renormalizados para que la fila sume } 1.
\]

\begin{table}[H]
\centering
\small
\renewcommand{\arraystretch}{1.4}
\begin{tabular}{@{}l l r@{}}
\toprule
\rowcolor{vialNavy!10}
\textbf{Tipo de evento} & \textbf{Ejemplo real \CDMX} & \textbf{Factor $f$} \\
\midrule
Protesta / marcha     & Marcha 9 de marzo, marchas CNTE & 1.70 \\
Festivo con cierres   & 15-sep (Grito), 16-sep (Desfile) & 1.80 \\
Cierre de línea Metro & Línea 1 (2021--22), Línea 12 (2022) & 1.55 \\
Evento masivo         & City Banamex, Foro Sol, Azteca & 1.40 \\
Día hábil base        & --- & 1.00 \\
Temporada baja        & Navidad, Año Nuevo & 0.60 \\
\bottomrule
\end{tabular}
\caption{Catálogo de perturbaciones contextuales (extracto). La función
         \texttt{seleccionar\_perturbacion(fecha, alcaldía)} devuelve
         el factor más severo aplicable para un datetime y alcaldía dados.}
\label{tab:perturbaciones}
\end{table}

\subsection{Simulación Monte Carlo vectorizada}

Con la matriz $P'$ vigente, el motor ejecuta $N = 10000$ trayectorias
en paralelo usando NumPy vectorizado (\texttt{src/simulation/monte\_carlo.py}).
Cada trayectoria simula los estados de tráfico segmento a segmento
desde el origen hasta el destino y acumula el tiempo de viaje total.
Los percentiles $P_{10}$, $P_{50}$, $P_{90}$ se extraen directamente
de la distribución empírica resultante, junto con el Índice de
Confiabilidad:

\[
IC = \frac{P_{90} - P_{10}}{P_{50}}
\]

Un IC $< 0.5$ (verde), $0.5$--$1.0$ (amarillo), $> 1.0$ (rojo) indica
si la ruta es confiable, variable o altamente incierta.

% ═══════════════════════ §4 AGENTE VialAI ══════════════════════
\section{El agente conversacional \VialAI}

El agente implementa un bucle de \emph{tool use} con Claude claude-sonnet-4-6
usando el SDK de Anthropic. El bucle sigue el patrón:

\begin{enumerate}
\item El usuario envía un mensaje (texto o voz transcrita).
\item El LLM identifica qué herramientas necesita y las llama.
\item Los resultados de las herramientas se inyectan en el contexto.
\item El LLM razona sobre los resultados y genera la respuesta final.
\end{enumerate}

\subsection{Las seis herramientas del agente}

\begin{table}[H]
\centering
\small
\renewcommand{\arraystretch}{1.4}
\begin{tabularx}{\linewidth}{@{}l X@{}}
\toprule
\rowcolor{vialNavy!10}
\textbf{Herramienta} & \textbf{Función} \\
\midrule
\texttt{predecir\_tiempo} & Ejecuta el motor Monte Carlo y devuelve
  $P_{10}$/$P_{50}$/$P_{90}$, distancia, velocidad media e IC. \\
\texttt{obtener\_trafico\_actual} & Consulta TomTom Traffic Stats para el
  estado actual de congestión en el corredor solicitado. \\
\texttt{obtener\_clima} & Consulta OpenWeatherMap para la condición
  climática vigente y ajusta el factor de velocidad. \\
\texttt{comparar\_rutas} & Ejecuta predicciones en paralelo para
  múltiples rutas y devuelve un ranking por IC. \\
\texttt{evaluar\_perturbacion} & Consulta el catálogo de perturbaciones
  activas y devuelve el factor $f$ aplicable. \\
\texttt{obtener\_lugares\_guardados} & Recupera los lugares personalizados
  del usuario (inyectados en el system prompt). \\
\bottomrule
\end{tabularx}
\caption{Las seis herramientas registradas en el agente \VialAI.}
\label{tab:herramientas}
\end{table}

\subsection{Interfaz de voz de tres capas}

La entrada y salida de voz siguen una arquitectura de tres capas con
degradación elegante (\emph{graceful degradation}):

\begin{enumerate}
\item \textbf{STT primario:} faster-whisper local (sin API externa,
sin costo, sin latencia de red).
\item \textbf{STT fallback:} Google Speech-to-Text vía API si el
modelo local falla.
\item \textbf{STT último recurso:} campo de texto manual en la UI.
\end{enumerate}

Para TTS (síntesis de voz):
\begin{enumerate}
\item \textbf{TTS premium:} voz \texttt{nova} de OpenAI (natural,
\textless 1 s de latencia).
\item \textbf{TTS local:} pyttsx3 (robótica pero funcional, sin
dependencias externas).
\item \textbf{Fallback silencioso:} el texto aparece en pantalla sin
reproducción de audio.
\end{enumerate}

La compresión semántica a \textless 30 palabras antes de sintetizar
está implementada en \texttt{voice\_io.py} usando Claude Haiku 4.5
(modelo más rápido y económico del stack).

% ═══════════════════════════════════════════════════════════════
%  PARTE II: CÓMO EL USUARIO USA EL SISTEMA
% ═══════════════════════════════════════════════════════════════

\newpage
\part*{Parte II — Cómo el usuario usa el sistema}
\addcontentsline{toc}{part}{Parte II: Cómo el usuario usa el sistema}

% ═══════════════════════ §5 INICIO ═════════════════════════════
\section{Inicio rápido}

Para ejecutar \VialAI\ en local (\textasciitilde 5 minutos):

\begin{lstlisting}[language=bash, caption={Instalación y ejecución de VialAI.}]
# 1. Clonar el repositorio
git clone https://github.com/caledelta/UrbanFlow_CDMX.git
cd UrbanFlow_CDMX

# 2. Crear entorno virtual e instalar dependencias
python -m venv .venv
source .venv/bin/activate      # Linux/Mac
.venv\Scripts\activate         # Windows
pip install -r requirements.txt

# 3. Configurar API keys
cp .env.example .env
# Editar .env y agregar: ANTHROPIC_API_KEY, TOMTOM_API_KEY,
#   OPENWEATHERMAP_API_KEY, GOOGLE_MAPS_API_KEY

# 4. (Opcional) Verificar que los tests pasan
pytest tests/ -q

# 5. Lanzar la aplicacion
streamlit run app/streamlit_app.py
\end{lstlisting}

La aplicación estará disponible en \texttt{http://localhost:8501} en el
navegador por defecto. Las API keys se leen exclusivamente desde
\texttt{.env}; nunca se escriben en el código.

% ═══════════════════════ §6 FLUJOS DE USO ══════════════════════
\section{Los tres flujos de uso}

La interfaz es \emph{single-page} con una barra lateral izquierda
(configuración de consulta) y un área principal (mapa + resultados).
El usuario puede interactuar de tres formas equivalentes:

\subsection{Flujo 1: predicción manual}

\begin{enumerate}
\item Seleccionar \textbf{Origen} y \textbf{Destino} desde los menús
desplegables o haciendo clic sobre el mapa.
\item Ajustar la \textbf{hora de salida} y la \textbf{fecha} (por
defecto: ahora).
\item Activar opcionalmente una \textbf{perturbación contextual} en
el selector (``Marcha'', ``Cierre Metro L12'', etc.).
\item Pulsar el botón \textcolor{vialTeal}{\textbf{Predecir trayecto}}.
\item Leer los resultados en la tarjeta central:
$P_{50}$ (mediana), $P_{10}$ (optimista), $P_{90}$ (adverso),
distancia y el Índice de Confiabilidad con su semáforo.
\end{enumerate}

\subsection{Flujo 2: chat con el agente}

\begin{enumerate}
\item Expandir el panel \textbf{Chat con VialAI} en la barra lateral.
\item Escribir la consulta en lenguaje natural.
\item El agente razona, llama a las herramientas necesarias (motor
Monte Carlo, tráfico actual, clima, lugares guardados) y responde
con la predicción integrada en un párrafo accionable.
\end{enumerate}

\begin{casebox}[Ejemplo de consulta al agente]
\textbf{Usuario:} \textit{``¿Vale la pena salir ahora para Santa Fe,
o espero 30 minutos? Estoy en Polanco, son las 8:10.''} \\[0.3em]
\textbf{VialAI:} \textit{``Ruta de 17.3~km. Mediana esperada: 27~min
(llegada 8:37). El IC es amarillo: en el 10\% adverso tardas hasta
44~min. Si tu cita es antes de las 9, sal ahora. Si tienes hasta las
9:30, esperar 30~min mejora el IC a verde.''} \\[0.2em]
\textit{[El agente llamó internamente a: \texttt{predecir\_tiempo},
\texttt{obtener\_trafico\_actual}, \texttt{obtener\_clima}]}
\end{casebox}

\subsection{Flujo 3: entrada por voz (manos libres)}

\begin{enumerate}
\item Pulsar el botón del micrófono \textbf{🎙} en la barra lateral.
\item Hablar la consulta (mismo lenguaje natural que en el chat).
\item Detener la grabación. El STT transcribe el audio localmente.
\item El agente procesa la consulta y sintetiza la respuesta en audio
(\textless 30 palabras, comprensible en una sola escucha).
\item Si el TTS falla, la respuesta aparece en texto como fallback.
\end{enumerate}

\begin{tipbox}[Diseño para el conductor en movimiento]
Los \textless 30 palabras de respuesta de voz no son arbitrarios:
un conductor en motocicleta debe comprender la respuesta en una
escucha sin apartar los ojos del tráfico. La compresión semántica
con Haiku preserva solo la información accionable: el tiempo esperado,
el nivel de confiabilidad y la recomendación concreta (``sal ahora''
o ``espera N minutos'').
\end{tipbox}

% ═══════════════════════ §7 FLUJO DE VIAJE ═════════════════════
\section{Flujo completo de un viaje}

Además de las predicciones puntuales, \VialAI\ soporta un flujo de
viaje completo que permite capturar el tiempo real y compararlo con
la predicción:

\begin{enumerate}
\item \textbf{Predecir}: el usuario solicita una predicción (cualquier
flujo de los tres anteriores).
\item \textbf{Iniciar viaje}: una vez en ruta, pulsar
\emph{``Iniciar viaje''}. El sistema registra el timestamp de inicio.
\item \textbf{Llegar}: al arribar, pulsar \emph{``Llegué''}. El sistema
calcula $\Delta = t_\text{real} - \hat{t}_{P50}$ y lo clasifica como
\emph{exacto} ($|\Delta| \leq 10\%$), \emph{cerca} ($\leq 25\%$) o
\emph{lejos} ($> 25\%$).
\item \textbf{Calificar}: aparece un rating de 1--5 estrellas opcional
que el usuario puede responder o posponer.
\item \textbf{Mis estadísticas}: el panel acumula el historial de
exactitud del sistema para ese usuario específico.
\end{enumerate}

Todos los datos se guardan \textbf{localmente} en
\texttt{data/telemetry/}. Ningún dato personal se transmite a
servidores externos.

% ═══════════════════════ §8 LUGARES GUARDADOS ══════════════════
\section{Lugares personalizados y marcadores en el mapa}

El usuario puede guardar sus lugares frecuentes (casa, trabajo,
clientes recurrentes) asignándoles un nombre y una categoría con
ícono contextual:

\begin{table}[H]
\centering
\small
\renewcommand{\arraystretch}{1.4}
\begin{tabular}{@{}l l l@{}}
\toprule
\rowcolor{vialNavy!10}
\textbf{Categoría} & \textbf{Ícono en el mapa} & \textbf{Ejemplo} \\
\midrule
Casa & 🏠 casa & ``Mi casa'' \\
Trabajo & 🏢 edificio & ``Oficina'' \\
Restaurante & 🍴 cubiertos & ``El Taquero de confianza'' \\
Gasolinera & ⛽ bomba & ``Pemex Ejército Nacional'' \\
Médico / Hospital & 🏥 hospital & ``Clínica 36 IMSS'' \\
Favorito genérico & ⭐ estrella & ``Punto de reunión'' \\
\bottomrule
\end{tabular}
\caption{Categorías de lugares guardados con sus íconos contextuales en el mapa Folium.}
\label{tab:lugares}
\end{table}

Los lugares guardados se inyectan automáticamente en el \emph{system
prompt} del agente, de modo que el usuario puede preguntar
\textit{``¿Cuánto tardo a mi trabajo?''} sin necesidad de especificar
la dirección cada vez.

% ═══════════════════════ §9 RECOMENDACIONES DE USO ═════════════
\section{Recomendaciones para el profesor / evaluador}

Para verificar las capacidades del sistema en una sesión corta,
se sugieren tres pruebas concretas:

\begin{enumerate}
\item \textbf{Predicción base} --- Origen: \textbf{Polanco},
Destino: \textbf{Santa Fe}, Hora: 08:15. Observar la banda
$P_{10}$/$P_{50}$/$P_{90}$ y el Índice de Confiabilidad.

\item \textbf{Activar una perturbación} --- En el panel lateral
activar \textbf{``Marcha 9 de marzo''} y recalcular. Notar cómo la
mediana sube y el IC pasa de amarillo a rojo.

\item \textbf{Chat con el agente} --- Preguntar: \textit{``¿Vale la
pena salir ahora para Santa Fe, o espero 30 minutos?''} Observar cómo
\VialAI\ razona y entrega una recomendación accionable.
\end{enumerate}

\begin{tipbox}[Si no se tienen API keys]
El sistema soporta \emph{modo offline}. Sin API keys externas,
el motor estocástico opera con los datos en caché local
(\texttt{data/}) y entrega predicciones igualmente válidas
para las rutas que ya han sido consultadas previamente.
El agente de voz y el motor Monte Carlo son totalmente locales
y no requieren credenciales externas.
\end{tipbox}

\end{document}
